\section{Observation graphs and physical priors}\label{sec:graphpriors}

\subsection{Coverage of microscopic vertices}
Let $\mathcal O$ be a nonempty reverse-closed set of individually resolved
microscopic transitions on a fixed, known vertex set $V$. Each observed
transition has a positive rate and known endpoints. The observation graph
has vertices $V$ and the undirected pairs underlying $\mathcal O$ as edges.
A vertex is covered if it is incident to an observed edge; an uncovered
vertex is fully hidden. Writing $E$ for the sum of observed jump matrices,
we retain $T=Q-E$, $R_I=e_{\mathrm{dest}(I)}^T$, and
$B_J=q_J e_{\mathrm{orig}(J)}$. The joint kernel is $Re^{Tt}B$.

\begin{proposition}[Recovery under full vertex coverage]\label{prop:coverage}
If every microscopic vertex is covered, two distinct positive Laplace
arguments determine the generator uniquely on the declared vertex set,
even when the observed rates are unknown. The inverse on admissible data
is rational and locally Lipschitz. It is therefore locally Lipschitz in
maximum row total variation and continuous under weak row-law convergence,
without a common rate cap. One positive argument suffices if all observed
rates are known.
\end{proposition}
\begin{proof}
For each vertex $i$, choose an observed mark $I_i$ ending at $i$ and an
observed mark $J_i$ beginning at $i$. Put $D=\diag(q_{J_i})$ and select
the square submatrix $F(\lambda)_{ij}=\widehat\Psi_{I_i,J_j}(\lambda)$.
Then
\begin{equation}
 F(\lambda)=(\lambda I-T)^{-1}D.
\end{equation}
For $M_k=F(\lambda_k)^{-1}$ and $\lambda_1\ne\lambda_2$,
\begin{equation}
 M_2-M_1=(\lambda_2-\lambda_1)D^{-1},\quad
 D=(\lambda_2-\lambda_1)(M_2-M_1)^{-1},\quad
 T=\lambda_1 I-DM_1.\label{eq:coverageinverse}
\end{equation}
For an additional observed mark $J$, let $f_J$ be its measured column
restricted to rows $I_i$. The identity
$(\lambda_1I-T)f_J=q_J e_{\mathrm{orig}(J)}$ recovers its rate.
Thus $Q=T+E$ is determined. The inverses have nonzero denominators at
every admissible source. They remain smooth in a neighborhood, and the
Laplace measurements are bounded continuous tests of the row laws.
When the observed rates are known, $D$ and $E$ are known and the last
formula in~\eqref{eq:coverageinverse} needs only one argument.
\end{proof}

At a complete source, Proposition~\ref{prop:coverage} also gives a locally
Lipschitz entropy inverse and a finite local upper ceiling. The same is
true at an irreducible sparse source if competitors have the same known
support. This condition requires fewer observations than all physical
edges: observing $1\leftrightarrow2$ and $2\leftrightarrow3$ covers a
complete three-state network while leaving $1\leftrightarrow3$ hidden.
One unobserved reciprocal pair is therefore not sufficient for unbounded
entropy inference.

The support condition remains consequential. Start from the unit tree
$1\leftrightarrow2\leftrightarrow3$ and observe both tree edges. Open
the unobserved pair at rates $q_{13}=e$ and
$q_{31}=\exp(-1/e^2)$, keeping all four old rates fixed. The generators
and their joint row laws converge to the tree, whereas the new pair
contributes $\sig\sim1/(3e)$. Full vertex coverage identifies every
generator in the sequence but does not make entropy continuous across
this newly opening support. This variant changes the trace by a
vanishing amount; a trace-preserving version needs unobserved rate slack.

At the other extreme, the existing hidden-pair construction applies
whenever two vertices are uncovered: its similarity fixes every observed
jump matrix, including when many edges are observed. Unknown support
then allows unbounded entropy in every joint-kernel neighborhood at fixed
dimension, observed rates and trace. These are two sufficient graph
criteria; they do not exhaust arbitrary observation graphs with exactly
one uncovered vertex. Nor do they extend directly to lumped states or
blurred transition labels. A macrostate event need not reset to a known
microscopic state, and consecutive waiting laws then need not determine
the entire observed path law~\cite{SkinnerDunkel2021PNAS,Bruno2005,FritzSeifert2026}.

\subsection{Kinetic and thermodynamic regularization}
Let
\begin{equation}
 \mathcal K(Q)=\sum_i\pi_i\sum_{j\ne i}q_{ij},\qquad
 a_{\max}(Q)=\max_{q_{ij}>0}\left|\log\frac{q_{ij}}{q_{ji}}\right|.
\end{equation}
Here $\mathcal K$ is the microscopic activity, in inverse-time units;
$a_{\max}$ is dimensionless. Stationary cancellation gives the elementary
bound $\sig\le\mathcal K a_{\max}$. A sharper established steady-state
bound is~\cite{NishiyamaHasegawa2023}
\begin{equation}
 \sig(Q)\le\mathcal K(Q)a_{\max}(Q)
           \tanh\!\left(\frac{a_{\max}(Q)}2\right).\label{eq:sharpprior}
\end{equation}
The right side is zero at $a_{\max}=0$. Thus an activity ceiling
$\mathcal K\le K_0$ and an affinity ceiling $a_{\max}\le a_0$ give a
uniform finite entropy bound. An observed-event activity ceiling is not
by itself a ceiling on microscopic activity. If all supported rates lie
between $m>0$ and $M$, then
$\mathcal K\le(N-1)M$ and $a_{\max}\le\log(M/m)$; a fixed trace
$-\tr Q=\tau$ also gives $\mathcal K\le\tau$.

A rate ceiling alone fails by Proposition~\ref{prop:tree}. Conversely, an
affinity ceiling alone does not control arbitrarily fast dissipative
cycles. A rate floor alone also fails. In fact even a floor together with
an activity ceiling is insufficient: take $q_{12}=c\ge1$ and all five
other rates of a complete three-state chain equal to one. Direct
stationary solution gives
\begin{equation}
 \pi=\left(\frac1{c+2},\frac{2c+1}{3(c+2)},\frac13\right),\quad
 \mathcal K=\frac{3(c+1)}{c+2}<3,\quad
 \sig=\frac{c-1}{3(c+2)}\log c\longrightarrow\infty.
\end{equation}
The fast rate acts at a rarely occupied state. These priors therefore
control different quantities and cannot be interchanged.

There is a separate continuity statement at the generator level. On a
fixed connected support with $m\le q_{ij}\le M$, $Q\mapsto\sig(Q)$ is
uniformly Lipschitz on the compact admissible class. With a fixed
affinity bound, entropy is continuous under entrywise convergence to an
irreducible generator even if reciprocal pairs disappear: contributions
from pairs positive at the limit are smooth, and each disappearing
directed term is bounded in magnitude by $a_0\pi_iq_{ij}\to0$.
Locally the corresponding bounded-ratio cone even gives Lipschitz
continuity, since the derivatives of $x\log(x/y)$ are bounded when
$x/y$ stays in a compact positive interval. These conclusions concern
the entropy functional. They do not imply generator identification or
good conditioning of a partially observed inverse, and a compact
parameter class alone need not make entropy single valued on an exact
fiber. If entropy is identifiable throughout such a compact class, its
induced function on the observed image is continuous by the compact
quotient argument; a quantitative inverse bound still needs additional
analysis.

\subsection{Physical interpretation and measurement limitations}
The theorems concern the stationary network entropy-production
functional of an even-state CTMC with one channel per ordered pair.
They require no specified reservoir realization. Interpretation as
physical dissipation requires correctly resolved thermodynamic states
and channels obeying local detailed balance~\cite{Esposito2012}.
For a channel with intrinsic state entropies $s_i$, the rate ratio may
take the form
\begin{equation}
 \log\frac{q_{ij}^{\alpha}}{q_{ji}^{\alpha}}
 =\Delta s_{\mathrm{res}}^{\alpha}+s_j-s_i.
\end{equation}
At stationarity the mean state-entropy increment vanishes, so total and
medium entropy-production rates agree with the usual environmental
sign convention; their trajectory quantities differ by a boundary term.
The generator alone does not determine a unique heat current or
reservoir decomposition. Aggregating parallel channels can conceal
additional physical dissipation even when all state transitions are seen.

In the sparse-boundary construction, the vanishing reverse rate entails
an unbounded edge rate affinity and hence an unbounded per-transition
thermodynamic bias under local detailed balance. A fixed affinity prior
excludes that limit. Rate ceilings and floors are kinetic assumptions;
their ratio constrains the thermodynamic bias. A cycle affinity is the
sum of edge rate affinities and is distinct from a bound on each edge.

Finally, timestamp bins, missed events, uncertain endpoints and blurred
labels require an explicit observation model. Finite timing resolution
with every event retained is different from detector dead time; the
former is the setting of the finite-resolution comparison in
Ref.~\cite{FritzErtelSeifert2025}. A known noisy processing of the same
records contracts total variation and cannot eliminate an existing
indistinguishability obstruction. Such processing may also remove the
reset structure needed for our reconstruction. For general coarse or
semi-Markov observations, the time reverse must be induced from the
microscopic experiment; a naive reversal can create spurious
irreversibility when coarse-graining and time reversal do not
commute~\cite{Blom2024,BiskerReply2024}.
